\documentclass[12pt,a4paper]{article}
% ============================================================
%  Structural Persistence Series — Shared Preamble
% ============================================================
\usepackage{fontspec}
\setmainfont{Hiragino Mincho ProN}
\setsansfont{Hiragino Sans}
\setmonofont{Menlo}

\usepackage{iftex}
\ifXeTeX
  \XeTeXlinebreaklocale "ja"
  \XeTeXlinebreakskip=0pt plus 1pt minus 0.1pt
\fi

\usepackage[margin=22mm]{geometry}
\usepackage{amsmath,amssymb,amsthm}
\usepackage{xcolor}
\usepackage{graphicx}
\usepackage{hyperref}
\usepackage{booktabs}
\usepackage{fancyhdr}
% enumitem not available in this TeX installation

% --- Colours ------------------------------------------------
\definecolor{PaperAccent}{HTML}{1F4E79}
\definecolor{PaperGray}{HTML}{51606F}

\hypersetup{
  colorlinks=true,
  linkcolor=PaperAccent,
  urlcolor=PaperAccent,
  citecolor=PaperAccent,
  pdflang=ja
}
\urlstyle{same}

% --- Typography ---------------------------------------------
\setlength{\parindent}{1.5em}
\setlength{\parskip}{0pt}
\setlength{\emergencystretch}{3em}
\linespread{1.08}
\setlength{\abovedisplayskip}{8pt plus 2pt minus 4pt}
\setlength{\belowdisplayskip}{8pt plus 2pt minus 4pt}
\setlength{\abovedisplayshortskip}{6pt plus 2pt minus 2pt}
\setlength{\belowdisplayshortskip}{6pt plus 2pt minus 2pt}

% --- Section label names ------------------------------------
\renewcommand{\abstractname}{要旨}

% --- Theorem-like environments ------------------------------
\theoremstyle{plain}
\newtheorem{proposition}{命題}[section]
\newtheorem{lemma}[proposition]{補題}
\newtheorem{corollary}[proposition]{系}

\theoremstyle{definition}
\newtheorem{definition}{定義}[section]
\newtheorem{assumption}{条件}[section]

\theoremstyle{remark}
\newtheorem*{remark}{注}

% --- Running header -----------------------------------------
\pagestyle{fancy}
\fancyhf{}
\renewcommand{\headrulewidth}{0.4pt}
\fancyhead[L]{\small\textcolor{PaperGray}{\nouppercase{\leftmark}}}
\fancyhead[R]{\small\textcolor{PaperGray}{\thepage}}
\fancypagestyle{plain}{%
  \fancyhf{}
  \renewcommand{\headrulewidth}{0pt}
  \fancyfoot[C]{\small\textcolor{PaperGray}{\thepage}}
}

% --- Title block macro --------------------------------------
\newcommand{\PaperTitleBlock}[3]{%
  % #1 = title, #2 = subtitle, #3 = date
  \thispagestyle{plain}%
  \begin{center}
  {\LARGE\bfseries #1\par}
  \vspace{0.4em}
  {\normalsize #2\par}
  \vspace{1.0em}
  {\large Akihito Sunagawa\par}
  \vspace{0.3em}
  {\normalsize #3\par}
  \end{center}
  \vspace{1.6em}
}

% Legacy 2-arg version (backward compat)
\newcommand{\PaperTitleBlockSimple}[2]{%
  \PaperTitleBlock{#1}{}{#2}
}


\begin{document}

\PaperTitleBlock
  {The Balance Principle of Structural Persistence}
  {Net consumption, recovery, and recovery-aware structural persistence}
  {April 29, 2026}

\begin{abstract}
The minimal form of structural persistence gives the closed shrinkage kernel

\[
S=Me^{-L},
\]

where \(L\) is cumulative structural consumption measured as logarithmic shrinkage of a feasible region.

This paper extends that kernel to systems with repair, recovery, redundancy, rollback, reconfiguration, and external support. At each time step, let \(d_{t}\) be structural consumption and \(r_{t}\) be recovery. Define net structural consumption by

\[
b_{t}=d_{t}-r_{t}.
\]

For a finite horizon \(n\),

\[
B_{n}=\sum_{t<n}b_{t}.
\]

Then the feasible-region mass satisfies

\[
m(V^{(n)})=m(V^{(0)})e^{-B_{n}}.
\]

With the resource-side term,

\[
S_{n}=M_{n} e^{-B_{n}}.
\]

This does not replace the minimal kernel. If \(r_{t}=0\), then \(B_{n}=L_{n}\). The central claim is that, in recovery-aware systems, structural persistence is governed by net consumption rather than consumption alone.
\end{abstract}

\section{Question}\label{question}

Structures are not only consumed. Software can roll back or add redundancy. Organizations can repair routines or redistribute authority. Reasoning systems can rescope, refresh dependencies, or use external notes. Learning systems can replay, separate adapters, or resynchronize along dependencies.

For such systems, the relevant quantity is not only how much structure has been consumed. It is how much has been consumed relative to how much has been recovered.

This difference is the balance principle of structural persistence. Here ``balance'' means accounting, not equilibrium. It means that consumption and recovery are read on the same logarithmic scale and then subtracted.

\section{Two-step update}\label{two-step-update}

Let \(X\) be a system, and let

\[
V^{(t)}\subseteq X
\]

be the feasible region compatible with maintaining the target structural condition at time \(t\). A finite measure \(m\) is fixed in advance. The target structural condition, measure, observation unit, and time horizon are also fixed in advance.

Within one observation unit, write the update as

\[
\begin{aligned}
V_{t}^- :&= K_{t}(V^{(t)}), \\
\qquad \\
V^{(t+1)} :&= R_{t}(V_{t}^-).
\end{aligned}
\]

The action \(K_{t}\) consumes feasible structural region. The action \(R_{t}\) recovers feasible structural region.

This two-step notation is an accounting representation within a fixed observation unit. If the actual order of physical interventions changes the endpoint \(V^{(t+1)}\), then the net quantity below changes as well.

All log-ratio comparisons are read on positive finite masses. Reaching zero is treated as a collapse boundary.

\section{Consumption, recovery, and net consumption}\label{consumption-recovery-and-net-consumption}

Define structural consumption by

\[
d_{t}=-\log\frac{m(V_{t}^-)}{m(V^{(t)})}.
\]

Define recovery by

\[
r_{t}=\log\frac{m(V^{(t+1)})}{m(V_{t}^-)}.
\]

Both quantities live on the same logarithmic scale.

Define net structural consumption by

\[
b_{t}=d_{t}-r_{t}.
\]

Then the intermediate set cancels:

\[
\begin{aligned}
b_{t} \\
&= \\
-\log\frac{m(V^{(t+1)})}{m(V^{(t)})}.
\end{aligned}
\]

Thus, after the endpoint is fixed, the local net change is read directly as a log-ratio between the starting and ending feasible regions.

\section{Cumulative balance and exponential kernel}\label{cumulative-balance-and-exponential-kernel}

Let

\[
B_{n}=\sum_{t=0}^{n-1}b_{t}.
\]

Then

\[
m(V^{(n)})=m(V^{(0)})e^{-B_{n}}.
\]

Proof. Since

\[
e^{-b_{t}}=\frac{m(V^{(t+1)})}{m(V^{(t)})},
\]

we have

\[
\begin{aligned}
e^{-B_{n}} \\
&= \\
\prod_{0\le t<n}e^{-b_{t}} \\
&= \\
\prod_{0\le t<n}\frac{m(V^{(t+1)})}{m(V^{(t)})} \\
&= \\
\frac{m(V^{(n)})}{m(V^{(0)})}.
\end{aligned}
\]

Multiplying by \(m(V^{(0)})\) gives the result.

With the resource-side term,

\[
S_{n}=M_{n} e^{-B_{n}}.
\]

When the horizon is clear, write

\[
S=Me^{-B}.
\]

Here \(S_{n}\) is the structural persistence potential for the target structural condition. Reaching \(S_{n}=0\), or a pre-fixed threshold \(S_{n}\le S_{c}\), means reaching the structural maintenance boundary for that condition. Observationally, the boundary may appear as collapse, functional failure, halt, phase transition, or structural reorganization.

The term \(M_{n}\) is not derived from this pathwise identity. It is the resource-side scalar representing the effective maintenance surplus available for the target structural condition. The structural maintenance boundary can be reached through the feasible-region route represented by cumulative net consumption \(B_{n}\), the resource-side route represented by \(M_{n}=0\), or both. This paper does not claim a universal predictive decomposition of \(M_{n}\).

\section{Relation to the minimal form}\label{relation-to-the-minimal-form}

In the minimal form, recovery is not made explicit. Equivalently, set

\[
r_{t}=0.
\]

Then

\[
\begin{aligned}
b_{t}&=d_{t}, \\
\qquad \\
B_{n}&=L_{n}.
\end{aligned}
\]

Therefore

\[
S_{n}=M_{ne}^{-B_{n}}
\]

reduces to

\[
S_{n}=M_{ne}^{-L_{n}}.
\]

The balance principle is thus an extension of the minimal kernel, not a competing kernel. Recovery changes the exponent from cumulative consumption \(L\) to cumulative net consumption \(B\).

\section{Three local accounting regimes}\label{three-local-accounting-regimes}

The signs of \(b_{t}\) give three local accounting regimes.

{\def\LTcaptype{table} % do not increment counter
\begin{longtable}[]{@{}ll@{}}
\toprule\noalign{}
Condition & Reading \\
\midrule\noalign{}
\endhead
\bottomrule\noalign{}
\endlastfoot
(b\_t\textgreater0) & consumption-dominant \\
(b\_t=0) & locally maintained \\
(b\_t\textless0) & recovery-dominant \\
\end{longtable}
}

These are not equilibrium claims. They are local accounting labels on a pre-fixed observation unit.

If two adjacent intervals are merged, then

\[
\begin{aligned}
b_{t:t+2} \\
&= \\
-\log\frac{m(V^{(t+2)})}{m(V^{(t)})} \\
&= \\
b_{t}+b_{t+1}.
\end{aligned}
\]

The cumulative quantity \(B_{n}\) is additively coherent under aggregation. The local regime labels, however, are grain-relative. A fine-grained path with consumption followed by recovery may become a locally maintained interval after aggregation.

Therefore, empirical claims about consumption, maintenance, recovery, collapse boundaries, or hitting times require the observation unit, stage sequence, and time horizon to be pre-fixed.

\section{Application and limitation}\label{application-and-limitation}

The kernel of this paper is pathwise algebraic. It states that, once the target structural condition, measure, observation unit, consumption, and recovery are fixed,

\[
m(V^{(n)})=m(V^{(0)})e^{-B_{n}}.
\]

In real domains, one must still decide what structure is being maintained, what measure compares feasible regions, and what observation or inference indicators approximate \(d_{t}\), \(r_{t}\), or \(B_{t}\). Discovery on the same data is not support. Support requires frozen mappings, frozen indicators, baselines, metrics, splits, and decision rules, followed by holdout, future-surface, fresh-archive, or outside-rerun validation.

The paper does not claim that every domain has a unique natural \(V,m,d_{t},r_{t}\). It does not claim that expectation-level tendencies imply high-probability hitting bounds without bounded increments, moment-generating functions, Chernoff/KL structure, margin conditions, or other additional assumptions. It does not claim universal-law closure.

\section{Conclusion}\label{conclusion}

The minimal kernel is

\[
S=Me^{-L}.
\]

The recovery-aware balance kernel is

\[
S=Me^{-B}.
\]

Here

\[
B_{n}=\sum_{t<n}(d_{t}-r_{t}).
\]

Recovery does not destroy the exponential form. It changes what enters the exponent. Structural persistence in open systems is therefore governed by cumulative net consumption: what has been consumed, what has been recovered, and what remains as net structural loss.
\end{document}
